% Section 4: Implementation
\section{Implementation: AWS Bedrock Demonstration}

\subsection{System Architecture}

Our implementation demonstrates prompt-based API design through a pizza ordering service using Amazon Nova Micro (\texttt{us.amazon.nova-micro-v1:0}), selected for optimal cost-performance balance \cite{aws_bedrock_2025}. Figure \ref{fig:architecture} illustrates the serverless architecture.

\begin{figure}[h]
\centering
\includegraphics[width=0.42\textwidth]{architecture}
\Description{AWS Serverless Architecture}
\caption{AWS Serverless Architecture.}
\label{fig:architecture}
\end{figure}

Request flow: client \gls{https} request $\rightarrow$ \gls{api} Gateway (auth, rate limiting) $\rightarrow$ Lambda (prompt construction) $\rightarrow$ Bedrock (LLM inference) $\rightarrow$ Lambda (validation, retry) $\rightarrow$ structured \gls{json} response.

\subsection{Prompt Template Design}

Nova Micro requires minimal formatting. Temperature set low (~0.3) for consistent \gls{json}:

\begin{lstlisting}[language=Python, basicstyle=\footnotesize\ttfamily]
system_prompt = """
You're a pizza ordering API. 
Extract pizza order details.
Then respond ONLY with valid JSON:
{
    "size": "small|medium|large", 
    "toppings": ["..."],
    "crust": "thin|thick|regular", 
    "quantity": number
}

User request: {user_prompt}
"""
\end{lstlisting}

This template explicitly constrains output structure. ``Respond ONLY with valid \gls{json}'' reduces hallucination \cite{sahoo_systematic_2025}, \cite{chen_prompt_2025}. Response handling addresses non-determinism through three error categories: \gls{json} parsing, schema validation, and service errors, each triggering retry, rule-based extraction, or user notification. Token accounting via \texttt{inputTokens} and \texttt{outputTokens} enables cost monitoring.

\subsection{Results and Analysis}

\textbf{Testing Methodology:} We designed a stratified test suite: (1) Standard orders (50): typical configurations; (2) Edge cases (25): unusual combinations, dietary restrictions; (3) Invalid inputs (15): nonsensical requests, prompt injection; (4) Ambiguous queries (10): underspecified orders.

\textbf{Results:} 98\% accuracy on standard orders, 88\% on edge cases, 93\% on invalid handling, 94\% on ambiguous queries. Overall first-pass: 94\%. Failures concentrate in edge cases. Latency: $847ms$ ($p50$), $1,340ms$ ($p95$), acceptable for interactive use \cite{aws_bedrock_2025}.

\textbf{Cost Analysis:} Table \ref{tab:cost_comparison} compares models for structured extraction:

\begin{table}[h]
\centering
\caption{Cost-Performance Comparison (1M orders)}
\label{tab:cost_comparison}
\small
\begin{tabular}{|l|c|c|c|l|}
\hline
\textbf{Model} & \textbf{Acc.} & \textbf{Lat.} & \textbf{Cost} & \textbf{Best For} \\
\hline
Nova Micro & 94\% & $847ms$ & \$5 & Cost-sensitive \\
Phi-3-mini & 93\% & $780ms$ & \$8 & Edge deploy \\
Gemini Flash & 95\% & $650ms$ & \$12 & Low latency \\
Claude Haiku & 96\% & $920ms$ & \$15 & High accuracy \\
\hline
\end{tabular}
\end{table}

Nova Micro achieves 83\% cost reduction versus Llama 3 8B \cite{aws_nova_2025} while maintaining competitive accuracy (94\% vs. 95\%), validating our "smallest sufficient model" principle for simple structured extraction tasks.

\textbf{Limitations \& Security:} Ambiguous requests require clarification through generic fallbacks. Multi-turn conversations need session management. Complex dietary restrictions occasionally produce invalid JSON requiring retry.

Security testing across five attack vectors confirmed layered defence effectiveness: instruction override (e.g., ``Ignore previous instructions...''), role manipulation, context injection, and token exhaustion were all \textbf{BLOCKED} via system prompt isolation, role constraints, input sanitization, and API Gateway limits; output format manipulation remained \textbf{PARTIALLY} vulnerable, requiring Bedrock Guardrails. This validates layered security: template design, input validation, Bedrock Guardrails, and API-level controls \cite{owasp_top_2024, aws_bedrock_2025}.
